\section{The boundary layer at a facet}\label{sec:boundary-layer}

The local law of Theorem~\ref{thm:local} holds on compact subsets of the open face. This
section follows the law of the occupation times from the inside of a face $F$ to one of its
facets $G=F\setminus\{i\}$, where the time spent in $i$ is small. We prove an exact formula
for the law near the facet, which is the run expansion with the runs in $i$ pulled out
(Lemmas~\ref{lem:layer-exact} and~\ref{lem:layer-exact-discrete}). It reduces the layer to
one statement about a weighted chain on $G$, a local law with a small complex parameter. We
prove that statement for every facet (Theorem~\ref{thm:L-facet}), in closed form when $G$ has one
or two states (Lemma~\ref{lem:L-two}), and the layer theorem (Theorem~\ref{thm:layer}) follows for
every face with at least two states. We compute the constants of the layer for every irreducible $Q$
(Proposition~\ref{prop:layer-constants}). Their exponent matches the rate of the local law for
every irreducible $Q$. Their prefactor matches the constant of the local law when $Q$ is
reversible (Proposition~\ref{prop:layer-prefactor}), and we do not prove the prefactor match
without reversibility. The near-vertex theorem of Section~\ref{sec:vertex}
is the pair-face case of this picture on the lattice (Section~\ref{sec:layer-vertex}).

\subsection{Setting}\label{sec:layer-setting}

Let $F$ be an irreducible face with $k=|F|\ge2$, let $i\in F$, and put $G=F\setminus\{i\}$.
We assume that $G$ satisfies (H), so $G$ is irreducible and $\mu(G)>0$. The start mass
$\mu_i$ may be zero. For $j\in G$ put
\[
c^-_j=q_{ji},\qquad c^+_j=q_{ij},
\]
the rates into $i$ and out of $i$. Both vectors are nonnegative and nonzero, since $F$ is
irreducible. For $z\in\mathbb C$ let
\[
Q_G(z)=Q_G+z\,c^-(c^+)^\top,\qquad \mu^{(z)}=\mu_G+z\mu_ic^+,\qquad e^{(z)}=\ones+zc^- ,
\]
and write $Q_G(z)=A(z)-\operatorname{diag}d(z)$ with $A(z)$ the off-diagonal part and
$d(z)=q_G-z\,c^-\circ c^+$. A jump $j\to j'$ of $Q_G(z)$ of rate $zc^-_jc^+_{j'}$ stands for a
visit to $i$ that is entered from $j$ and left to $j'$.

\emph{Occupation coordinates.} Let $E_F$ be the event of Lemma~\ref{lem:continuum}. On $E_F$
let $s\in(0,\infty)^G$ be the times spent in the states of $G$, $S=\sum_{j\in G}s_j$, and
$x=T-S$ the time spent in $i$. We write
\[
f_F(s,x)=e^{-\langle q_G,s\rangle-q_ix}H_F(s,x),\qquad f_G(s)=e^{-\langle q_G,s\rangle}H_G(s),
\]
where $H_F$ and $H_G$ are the sums $H$ of Lemma~\ref{lem:continuum} for the faces $F$ and $G$.
By that lemma $f_F$ is the density of the occupation times on $E_F$ in any $k-1$ of the
coordinates, and $f_G(s)$ is the same density for $G$ at total time $S$. In this
notation Theorem~\ref{thm:local} for the face $G$ reads
\begin{equation}\label{eq:layer-local-G}
f_G(S\beta)=S^{-(k-2)/2}C_G(\beta)e^{-SJ_G(\beta)}\bigl(1+O_{\mathcal K}(1/S)\bigr)
\end{equation}
uniformly for $\beta$ in a compact $\mathcal K\subset\Delta_G^\circ$ and $S\ge1$, with
$C_G(\beta)=\mu_j$ and $J_G=q_j$ when $G=\{j\}$. The profile on the facet is
$\beta=s/S\in\Delta_G^\circ$.

\emph{The weighted chain.} For real $z\ge0$ the matrix $Q_G(z)$ is Metzler and irreducible.
We let $J_G(\beta,z)$ be its rate, the supremum in Proposition~\ref{prop:rate} with $Q_G(z)$
in place of $Q_F$, and $C_G(\beta,z)$ its local-law constant, which is
Definition~\ref{def:tilt} for the off-diagonal part $A(z)$ with the start $\mu^{(z)}$ in place
of $\mu$ and the end vector $e^{(z)}$ in place of $\ones$, namely
\[
C_G(\beta,z)=\frac{(\mu^{(z)}\cdot r_\beta(z))(l_\beta(z)\cdot e^{(z)})}{(2\pi)^{(k-2)/2}\sqrt{\det\Sigma_\beta(z)}} .
\]
So $J_G(\beta,0)=J_G(\beta)$ and $C_G(\beta,0)=C_G(\beta)$.

\emph{The facet constants.} Let $r_\beta=r_\beta(0)$ and $l_\beta=\beta/r_\beta$ be the vectors
of Definition~\ref{def:tilt} for the face $G$, so $l_\beta\cdot r_\beta=1$. Put
\[
a_\beta=c^-\cdot l_\beta,\qquad b_\beta=c^+\cdot r_\beta,\qquad \omega_\beta=a_\beta b_\beta,\qquad
\pi^{\rm in}_{\beta,j}=\frac{c^-_jl_{\beta,j}}{a_\beta},\qquad \pi^{\rm out}_{\beta,j}=\frac{c^+_jr_{\beta,j}}{b_\beta}.
\]
Only $\omega_\beta$ is intrinsic, since $a_\beta$ and $b_\beta$ move in opposite ways when
$r_\beta$ is scaled. The law $\pi^{\rm in}_\beta$ is the law of the state from which the
tilted chain enters $i$, and $\pi^{\rm out}_\beta$ the law of the state into which it leaves
$i$. Let $\mathcal L_\beta$ be the Laplacian on $G$ with the conductances
$\kappa_{jj'}=l_{\beta,j}q_{jj'}r_{\beta,j'}+l_{\beta,j'}q_{j'j}r_{\beta,j}$ for $j\ne j'$,
which are those of the additive reversibilization of the tilted generator
$\widehat Q_\beta$. Its kernel is $\R\ones$, since $G$ is irreducible. Let
$\mathcal L_\beta^+$ be its pseudo-inverse and
\begin{equation}\label{eq:layer-gamma}
\gamma_\beta=q_i+\frac{\omega_\beta}2\,(\pi^{\rm out}_\beta-\pi^{\rm in}_\beta)^\top\mathcal L_\beta^+(\pi^{\rm out}_\beta-\pi^{\rm in}_\beta).
\end{equation}
The quadratic form is the effective resistance of the network $\kappa$ between the laws
$\pi^{\rm in}_\beta$ and $\pi^{\rm out}_\beta$. So $\gamma_\beta\ge q_i$, with equality exactly
when the two laws agree. Finally put
\[
\zeta=2\sqrt{\omega_\beta Sx},\qquad z_*=\sqrt{\frac x{\omega_\beta S}},\qquad
\mathbf B_\beta(x,S)=e^{-\gamma_\beta x}\sqrt{\frac{\omega_\beta S}x}\,I_1(\zeta),
\]
with the modified Bessel function $I_1$. We call $\mathbf B_\beta$ the Bessel kernel.

\subsection{The exact representation}\label{sec:layer-exact}

\begin{lemma}[the runs in $i$ pulled out]\label{lem:layer-exact}
For $z\in\mathbb C$ and $s\in(0,\infty)^G$ let
\begin{gather*}
f^{(z)}_G(s)=e^{-\langle d(z),s\rangle}\sum_{m\in\Z_{\ge1}^G}W_z(m)\prod_{j\in G}\frac{s_j^{m_j-1}}{(m_j-1)!},\\
W_z(m)=\sum_w\mu^{(z)}_{w_1}\prod_{l<|w|}A(z)_{w_lw_{l+1}}\,e^{(z)}_{w_{|w|}},
\end{gather*}
the inner sum over the words $w$ in $G$ with no two equal neighbours and $m_j$ letters $j$.
Then $z\mapsto f^{(z)}_G(s)$ is entire, $f^{(0)}_G=f_G$, and for $x>0$ and every $\rho>0$
\begin{align}
f_F(s,x)&=e^{-q_ix}\sum_{K\ge1}\frac{x^{K-1}}{(K-1)!}\,[z^K]f^{(z)}_G(s)\label{eq:layer-series}\\
&=e^{-q_ix}\,\frac1{2\pi\mathrm i}\oint_{|z|=\rho}e^{x/z}f^{(z)}_G(s)\,\frac{dz}{z^2}.\label{eq:layer-contour}
\end{align}
\end{lemma}

So $f^{(z)}_G$ is the killed density on the facet of the chain in which each visit to $i$ is
replaced by an instantaneous jump $j\to j'$ of rate $zc^-_jc^+_{j'}$, and $z$ counts the
visits. The $K$ visits have total length $x$, with density $q_i^Kx^{K-1}e^{-q_ix}/(K-1)!$,
and each return carries a factor $c^+_{j'}/q_i$. So the powers of $q_i$ cancel.

\begin{proof}
Fix a word $c$ in $F$ with no equal neighbours, every letter of $F$, and $K=m_i(c)\ge1$
letters $i$. Its term in $H_F(s,x)$ is $\mu_{c_1}\prod q_{c_lc_{l+1}}$ times
$x^{K-1}/(K-1)!$ times $\prod_{j\in G}s_j^{m_j-1}/(m_j-1)!$. The letters $i$ are isolated in
$c$. We delete them and mark the places where they were. The result is a marked word $c'$ in
$G$, in which two equal neighbours occur exactly where $c$ had $j\,i\,j$. A mark between
letters $j$ and $j'$ carries $q_{ji}q_{ij'}=c^-_jc^+_{j'}$. A mark before the first letter $j'$
carries $\mu_ic^+_{j'}$, and a mark after the last letter $j$ carries $c^-_j$. Two neighbours
without a mark are different and carry $q_{jj'}$. So the weight of $c$ is the product of these
weights. Between two different letters of $c'$ there is at most one mark, and between two
equal letters exactly one.

We group the marked words by their reduced word $w$, obtained by merging each block of equal
letters of $c'$ into one letter, and by the numbers $n_1,n_2,\dots$ of marks inside the
blocks. If $j$ forms $\varrho$ blocks with $n$ inner marks in all, it has $\varrho+n$ letters
in $c'$, and there are $\binom{n+\varrho-1}{\varrho-1}$ ways to spread the marks. Since
$\binom{n+\varrho-1}{\varrho-1}/(n+\varrho-1)!=1/(n!(\varrho-1)!)$,
\[
\sum_{n\ge0}\binom{n+\varrho-1}{\varrho-1}\frac{y^ns^{n+\varrho-1}}{(n+\varrho-1)!}=e^{ys}\frac{s^{\varrho-1}}{(\varrho-1)!}.
\]
With $y=zc^-_jc^+_j$, one power of $z$ for each mark, the inner marks give the factor
$e^{zc^-_jc^+_js_j}$, which is the diagonal part of $d(z)$. The marks between different
letters give $A(z)$, and the first and last marks give $\mu^{(z)}$ and $e^{(z)}$. Collecting
the powers of $z$ gives~\eqref{eq:layer-series}. The entries of $A(z)$, $\mu^{(z)}$ and
$e^{(z)}$ are at most $q_{\max}(1+|z|q_{\max})$, $1+|z|q_{\max}$ and $1+|z|q_{\max}$ in
modulus. So the bound $W(m)\le k^Mq_{\max}^{M-1}$ of Section~\ref{sec:model} holds for $W_z$
with $q_{\max}$ replaced by $q_{\max}(1+|z|q_{\max})$, up to a factor $(1+|z|q_{\max})^2$, and
the series for $f^{(z)}_G(s)$ converges locally uniformly in $z$. Hence it is entire. For
$z=0$ no word carries a mark, and $f^{(0)}_G$ is the sum of Lemma~\ref{lem:continuum} for $G$.
Finally we expand $e^{x/z}=\sum_{n\ge0}x^nz^{-n}/n!$ and integrate term by term on the circle.
This gives $\sum_{n\ge0}(x^n/n!)[z^{n+1}]f^{(z)}_G(s)$, which is~\eqref{eq:layer-series}.
\end{proof}

The discrete law has the same structure. The only change is that the visits to $i$ come in
blocks of steps, whose lengths form a composition of $n_i$.

\begin{lemma}[the discrete form]\label{lem:layer-exact-discrete}
Let $hq_{\max}\le1$, $\supp n=F$, $N'=N-n_i$, and let $\hat n$ be $n$ restricted to $G$. For
$z\in\mathbb C$ let
\[
P^{(z)}_{N'}(\hat n)=\sum_{x}\mu^{h,(z)}_{x_1}\prod_{t<N'}\bigl(I+hQ_G+zh^2c^-(c^+)^\top\bigr)_{x_tx_{t+1}}e^{h,(z)}_{x_{N'}},
\]
the sum over the words $x_1\cdots x_{N'}$ in $G$ with composition $\hat n$, where
$\mu^{h,(z)}=\mu_G+zh\mu_ic^+$ and $e^{h,(z)}=\ones+zhc^-$. Then $P^{(0)}_{N'}(\hat n)$ is the
probability of the composition $\hat n$ for the chain run for $N'$ steps, and
\[
P_N(n)=\sum_{K\ge1}\binom{n_i-1}{K-1}(1-hq_i)^{n_i-K}\,[z^K]P^{(z)}_{N'}(\hat n).
\]
\end{lemma}

\begin{proof}
We delete the maximal blocks of $i$ from a path with composition $n$. What is left is a path
$x_1\cdots x_{N'}$ in $G$ with composition $\hat n$. A block of length $\ell$ between
$x_t=j$ and $x_{t+1}=j'$ replaces the factor $(I+hQ)_{jj'}$ by
$hq_{ji}(1-hq_i)^{\ell-1}hq_{ij'}$. A block at the start replaces $\mu_{x_1}$ by
$\mu_i(1-hq_i)^{\ell-1}hq_{ix_1}$, and a block at the end adds the factor
$hq_{x_{N'}i}(1-hq_i)^{\ell-1}$. We count the blocks by $z$. If there are $K$ blocks, the
factors $(1-hq_i)^{\ell-1}$ give $(1-hq_i)^{n_i-K}$, and the block lengths form one of the
$\binom{n_i-1}{K-1}$ compositions of $n_i$ into $K$ positive parts.
\end{proof}

\subsection{Pair faces}\label{sec:layer-pair}

For $k=2$ the facet is one state and the weighted chain has no words to sum. The contour
formula then gives the law in closed form.

\begin{proposition}[pair faces]\label{prop:pair-face}
Let $F=\{j,i\}$ with $q_{ji},q_{ij}>0$ and $\mu_j>0$, and put $\omega=q_{ji}q_{ij}$ and
$\zeta=2\sqrt{\omega sx}$. For all $s,x>0$
\[
f_F(s,x)=\mu_je^{-q_js-q_ix}\Bigl[\sqrt{\frac{\omega s}x}\,I_1(\zeta)+\Bigl(q_{ji}+\frac{\mu_i}{\mu_j}q_{ij}\Bigr)I_0(\zeta)
+\frac{\mu_i}{\mu_j}\sqrt{\frac{\omega x}s}\,I_1(\zeta)\Bigr].
\]
Here $\omega_\beta=\omega$, $\gamma_\beta=q_i$ and $f_G(s)=\mu_je^{-q_js}$, and for all
$s,x>0$
\[
0\le\frac{f_F(s,x)}{f_G(s)\mathbf B_\beta(x,s)}-1\le\Bigl(q_{ji}+\frac{\mu_i}{\mu_j}q_{ij}\Bigr)\Bigl(\frac1{\omega s}+\sqrt{\frac x{\omega s}}\Bigr)+\frac{\mu_i}{\mu_j}\,\frac xs .
\]
\end{proposition}

This is the facet version of the formula of Pedler for the occupation times of a
two-state chain~\cite{pedler1971}, and Kovchegov, Meredith and Nir give Bessel densities of
the same kind~\cite{kmn2010}. What is new is the reading as a layer with an explicit
two-sided bound. For $x\le T/2$ the bound is $O(1/T+\sqrt{x/T})$, and the term
$\sqrt{x/T}$ is attained, since the middle term of the formula has exactly this size.

\begin{proof}
For $G=\{j\}$ the only word is $j$, so
$f^{(z)}_G(s)=(\mu_j+z\mu_iq_{ij})(1+zq_{ji})e^{-(q_j-z\omega)s}$. With $z=z_*t$ and
$z_*=\sqrt{x/(\omega s)}$ we have $x/z+\omega sz=\frac\zeta2(t+1/t)$, and the generating
function $e^{(\zeta/2)(t+1/t)}=\sum_{n\in\Z}I_n(\zeta)t^n$~\cite[Eq.~10.35.1]{dlmf} gives
\begin{equation}\label{eq:layer-bessel-moments}
\frac1{2\pi\mathrm i}\oint e^{x/z+\omega sz}z^{n-2}\,dz=z_*^{n-1}I_{n-1}(\zeta)\qquad(n\in\Z).
\end{equation}
We put this into~\eqref{eq:layer-contour} for $n=0,1,2$ and use $I_{-1}=I_1$. This gives the
formula. Each term is nonnegative, which gives the lower bound. For the upper bound we divide
the last two terms by the first. This gives
$(q_{ji}+\frac{\mu_i}{\mu_j}q_{ij})z_*I_0(\zeta)/I_1(\zeta)+\frac{\mu_i}{\mu_j}x/s$. The lower
bound $I_1(\zeta)/I_0(\zeta)\ge\zeta/(1+\sqrt{1+\zeta^2})$ of Amos~\cite{amos1974} gives
$I_0/I_1\le(1+\sqrt{1+\zeta^2})/\zeta\le1+2/\zeta$, and $z_*(1+2/\zeta)=\sqrt{x/(\omega s)}+1/(\omega s)$.
For a pair $\beta=(1)$, the vectors of the facet are $r_\beta=l_\beta=1$, so
$\omega_\beta=\omega$ and $\pi^{\rm in}_\beta=\pi^{\rm out}_\beta$, which gives
$\gamma_\beta=q_i$.
\end{proof}

\subsection{The constants for every irreducible chain}\label{sec:layer-constants}

The rate on $F$ is homogeneous of degree one in the occupation times, and we write
$J_F(s,x)=TJ_F((s,x)/T)$. For $u\in(0,\infty)^G$ put
\[
\Psi_\beta(u)=\sum_{j\in G}\beta_j\Bigl(q_j-\sum_{j'\in G\setminus\{j\}}q_{jj'}\frac{u_{j'}}{u_j}\Bigr),\qquad
P_\beta(u)=(c^+\cdot u)\sum_{j\in G}\frac{c^-_j\beta_j}{u_j}.
\]
Since $\sum_j\beta_j(-(Q_G(z)u)_j/u_j)=\Psi_\beta(u)-zP_\beta(u)$, the rate of the weighted
chain is $J_G(\beta,z)=\sup_u[\Psi_\beta(u)-zP_\beta(u)]$ for real $z\ge0$.

\begin{proposition}[the facet expansion of the rate]\label{prop:layer-constants}
Let $F$ be irreducible, $i\in F$, and $G=F\setminus\{i\}$ irreducible.
\begin{enumerate}[(i)]
\item For $s=S\beta$ with $\beta\in\Delta_G^\circ$ and $x>0$,
\[
J_F(s,x)=q_ix+\sup_{z>0}\bigl[S\,J_G(\beta,z)-x/z\bigr].
\]
\item The map $z\mapsto J_G(\beta,z)$ is real-analytic near $z=0$, jointly with $\beta$, and
\[
J_G(\beta,0)=J_G(\beta),\qquad \partial_zJ_G(\beta,0)=-\omega_\beta,\qquad
\partial_z^2J_G(\beta,0)=2\omega_\beta(\gamma_\beta-q_i).
\]
\item There are $\delta>0$ and $C$ such that, uniformly for $\beta$ in a compact subset of
$\Delta_G^\circ$ and $0<x\le\delta S$,
\[
\Bigl|J_F(s,x)-\bigl(S\,J_G(\beta)+q_ix-2\sqrt{\omega_\beta Sx}+(\gamma_\beta-q_i)x\bigr)\Bigr|\le C\,x^{3/2}S^{-1/2}.
\]
\item If $\nu_jq_{jj'}=\nu_{j'}q_{j'j}$ on $F$ for some $\nu\in(0,\infty)^F$, then
$\pi^{\rm in}_\beta=\pi^{\rm out}_\beta$, $\gamma_\beta=q_i$, $J_G(\beta,z)=J_G(\beta)-z\omega_\beta$
for all real $z\ge0$, $\omega_\beta=\bigl(\sum_{j\in G}\sqrt{q_{ij}q_{ji}\beta_j}\bigr)^2$, and
for all $s=S\beta$ and $x>0$
\[
J_F(s,x)=S\,J_G(\beta)+q_ix-2\sqrt{\omega_\beta Sx}.
\]
\end{enumerate}
\end{proposition}

So in the visit fraction $\alpha_i$ the rate has a square-root singularity
$-2\sqrt{\omega_\beta\alpha_i}$ at the facet, and the time in $i$ is charged at the rate
$\gamma_\beta$, not $q_i$.

\begin{proof}
(i) In the supremum that defines $J_F(s,x)$ we write $u=(u_G,\epsilon)$ with
$\epsilon=u_i$. The terms with $\epsilon$ are
$-\epsilon\sum_js_jc^-_j/u_j-(x/\epsilon)(c^+\cdot u_G)$, and both coefficients are positive.
Their supremum over $\epsilon>0$ is $-2\sqrt{xS\,P_\beta(u_G)}$. So
$J_F(s,x)=q_ix+\sup_u[S\Psi_\beta(u)-2\sqrt{xSP_\beta(u)}]$. Now
$2\sqrt{xSP}=\inf_{z>0}[SzP+x/z]$, and we exchange the two suprema.

(ii) Write $u=e^\psi$. For real $z\ge0$ the function $\psi\mapsto\Psi_\beta(e^\psi)-zP_\beta(e^\psi)$
is concave, since $P_\beta(e^\psi)$ is a positive combination of the convex functions
$e^{\psi_m-\psi_j}$, and it does not change when a constant is added to $\psi$. At $z=0$ its
maximizer is $\psi_0=\log r_\beta$, unique up to a constant (Proposition~\ref{prop:rate}(b)).
The Hessian of $\sum_{j\ne j'}\beta_jq_{jj'}e^{\psi_{j'}-\psi_j}$ at $\psi_0$ is the Laplacian
with the conductances $\kappa_{jj'}$, because $\beta_jq_{jj'}r_{\beta,j'}/r_{\beta,j}=l_{\beta,j}q_{jj'}r_{\beta,j'}$.
So the Hessian of $\Psi_\beta$ at $\psi_0$ is $-\mathcal L_\beta$, which is negative definite on
the complement of the constants. By the implicit function theorem on that complement, the
critical point $\psi(z)$ is real-analytic in $(\beta,z)$ near $(\beta,0)$, and for real $z\ge0$
it is the maximizer by concavity. So $J_G(\beta,z)=\Psi_\beta(\psi(z))-zP_\beta(\psi(z))$ is
real-analytic. The envelope theorem gives $\partial_zJ_G=-P_\beta(\psi(z))$, so
$\partial_zJ_G(\beta,0)=-(c^+\cdot r_\beta)(c^-\cdot l_\beta)=-\omega_\beta$. Differentiating the
critical equation gives $-\mathcal L_\beta\psi'(0)=\nabla P$, with $\nabla P$ the gradient of
$P_\beta(e^\psi)$ at $\psi_0$. Its entries are
\[
\partial_{\psi_m}P=c^+_mr_{\beta,m}\,(c^-\cdot l_\beta)-(c^+\cdot r_\beta)\,c^-_ml_{\beta,m}=\omega_\beta(\pi^{\rm out}_{\beta,m}-\pi^{\rm in}_{\beta,m}),
\]
which sum to zero. So $\partial_z^2J_G(\beta,0)=-\nabla P\cdot\psi'(0)=\nabla P^\top\mathcal L_\beta^+\nabla P$,
and this is $2\omega_\beta(\gamma_\beta-q_i)$ by~\eqref{eq:layer-gamma}.

(iii) By the proof of (i), $J_F(s,x)-q_ix=S\sup_\psi[\Psi_\beta(\psi)-2\varepsilon P_\beta(\psi)^{1/2}]$
with $\varepsilon=\sqrt{x/S}$. The maximizer lies in a bounded set of $\psi$ modulo constants,
uniformly for $\beta$ in a compact set, as in the proof of Proposition~\ref{prop:rate}(b).
Comparing its value with the value at $\psi_0$ gives
$\Psi_\beta(\psi_0)-\Psi_\beta(\psi^*)\le2\varepsilon\sup P_\beta^{1/2}$ on that set, and the
strict concavity near $\psi_0$ then gives $\psi^*-\psi_0=O(\varepsilon^{1/2})$. The critical
equation $\nabla\Psi_\beta(\psi^*)=\varepsilon\nabla P_\beta(\psi^*)/\sqrt{P_\beta(\psi^*)}$ has a
right side of size $O(\varepsilon)$, and $\nabla^2\Psi_\beta$ is invertible near $\psi_0$ on the
complement of the constants, so $\psi^*-\psi_0=O(\varepsilon)$. Now we expand to
second order around $\psi_0$. With $g=\nabla(P^{1/2})=\nabla P/(2\sqrt{\omega_\beta})$, the maximizer
is $\psi^*=\psi_0-2\varepsilon\mathcal L_\beta^+g+O(\varepsilon^2)$, and the value is
$J_G(\beta)-2\varepsilon\sqrt{\omega_\beta}+2\varepsilon^2g^\top\mathcal L_\beta^+g+O(\varepsilon^3)$.
Here $2g^\top\mathcal L_\beta^+g=\nabla P^\top\mathcal L_\beta^+\nabla P/(2\omega_\beta)=\gamma_\beta-q_i$
by (ii). Multiplying by $S$ gives (iii), since $S\varepsilon^3=x^{3/2}S^{-1/2}$.

(iv) By Proposition~\ref{prop:constant}(b) for the face $G$ we may take
$r_{\beta,j}=\sqrt{\beta_j/\nu_j}$ and $l_{\beta,j}=\sqrt{\beta_j\nu_j}$. Then
$c^-_jl_{\beta,j}=q_{ji}\sqrt{\beta_j\nu_j}=\nu_iq_{ij}\sqrt{\beta_j/\nu_j}=\nu_ic^+_jr_{\beta,j}$, so
$\pi^{\rm in}_\beta=\pi^{\rm out}_\beta$ and $\gamma_\beta=q_i$. More is true. For every real
$z\ge0$ the matrix $Q_G(z)$ is reversible with respect to $\nu_G$, since
$\nu_jc^-_jc^+_{j'}=\nu_iq_{ij}q_{ij'}$ is symmetric in $j,j'$. So by
Proposition~\ref{prop:constant}(b) its optimal vector is $r_\beta$ for every $z$, and
$J_G(\beta,z)=\Psi_\beta(r_\beta)-zP_\beta(r_\beta)=J_G(\beta)-z\omega_\beta$. Next
$\omega_\beta=\nu_i\bigl(\sum_jq_{ij}\sqrt{\beta_j/\nu_j}\bigr)^2=\bigl(\sum_j\sqrt{q_{ij}q_{ji}\beta_j}\bigr)^2$,
using $q_{ji}=q_{ij}\nu_i/\nu_j$. Finally (i) gives
$\sup_{z>0}[-Sz\omega_\beta-x/z]=-2\sqrt{\omega_\beta Sx}$.
\end{proof}

For $1/T\ll x\ll T$ the Bessel kernel is
$(4\pi)^{-1/2}(\omega_\beta S)^{1/4}x^{-3/4}\exp(2\sqrt{\omega_\beta Sx}-\gamma_\beta x)(1+O(1/\zeta))$,
by $I_1(\zeta)=e^\zeta(2\pi\zeta)^{-1/2}(1+O(1/\zeta))$~\cite[Eq.~10.40.1]{dlmf}. So
Proposition~\ref{prop:layer-constants}(iii) says that $f_G(s)\mathbf B_\beta(x,S)$ and the local
law of $F$, which in occupation coordinates is $T^{-(k-1)/2}C_F(\alpha)e^{-TJ_F(\alpha)}$, have
the same exponent up to $O(x^{3/2}S^{-1/2})$. For their powers of $x$ and their constants to
agree as well we need $C_F((1-\epsilon)\beta,\epsilon)\sim(4\pi)^{-1/2}\omega_\beta^{1/4}\epsilon^{-3/4}C_G(\beta)$
as $\epsilon\to0$. The next proposition proves this for reversible $Q$.

\subsection{Prefactor matching}\label{sec:layer-prefactor}

\begin{proposition}[prefactor matching]\label{prop:layer-prefactor}
Let $Q$ be reversible on $F$ with respect to $\nu$, and let $G$ satisfy (H), with $|G|\ge1$.
Uniformly for $\beta$ in a compact subset of $\Delta_G^\circ$, as $\epsilon\to0$,
\[
C_F((1-\epsilon)\beta,\epsilon)=(4\pi)^{-1/2}\,\omega_\beta^{1/4}\,\epsilon^{-3/4}\,C_G(\beta)\bigl(1+O(\sqrt\epsilon)\bigr),
\]
with $C_G(\beta)=\mu_j$ when $G=\{j\}$.
\end{proposition}

\begin{proof}
We use Proposition~\ref{prop:constant}(b) at $\alpha=((1-\epsilon)\beta,\epsilon)$. On $G$ the
vectors $r_\alpha=\sqrt{\alpha/\nu}$ and $l_\alpha=\sqrt{\alpha\nu}$ are $\sqrt{1-\epsilon}$ times
$r_\beta$ and $l_\beta$, and at $i$ they are $O(\sqrt\epsilon)$. So
$(\mu\cdot r_\alpha)(l_\alpha\cdot\ones)=(\mu_G\cdot r_\beta)(l_\beta\cdot\ones)(1+O(\sqrt\epsilon))$,
where $\mu_G\cdot r_\beta>0$ by (H). Next $\prod_{j\in F}\alpha_j=\epsilon(1-\epsilon)^{k-1}\prod_{j\in G}\beta_j$.
The conductances are $c_{jj'}(\alpha)=(1-\epsilon)c_{jj'}(\beta)$ on $G$, and at $i$ they are
$c_{ij}(\alpha)=\sqrt{\epsilon(1-\epsilon)}\sqrt{q_{ij}q_{ji}\beta_j}$. A spanning tree of $F$
with $m$ edges at $i$ has weight $O(\epsilon^{m/2})$. The trees with one edge at $i$ are a
spanning tree of $G$ and one edge $\{i,j\}$. So
$\tree(c(\alpha))=\sqrt\epsilon\,\sqrt{\omega_\beta}\,\tree(c_G(\beta))(1+O(\sqrt\epsilon))$, by
$\sum_j\sqrt{q_{ij}q_{ji}\beta_j}=\sqrt{\omega_\beta}$ from
Proposition~\ref{prop:layer-constants}(iv). We put these into the formula for $C_F$ and divide
by the same formula for $C_G$, which has one factor $(4\pi)^{1/2}$ fewer. The powers of
$\epsilon$ are $\epsilon^{1/4}$ from the tree and $\epsilon^{-1}$ from the product, which gives
$\epsilon^{-3/4}$. For $G=\{j\}$ we use $C_{\{j\}}=\mu_j$ and $\tree=1$ on one vertex.
\end{proof}

\emph{The first amplitude correction.} Put $c_1(\beta)=\partial_z\log C_G(\beta,z)$ at $z=0$.
For reversible $Q$,
\begin{equation}\label{eq:layer-c1}
c_1(\beta)=\frac{\mu_ib_\beta}{\mu_G\cdot r_\beta}+\frac{a_\beta}{l_\beta\cdot\ones}+\frac{\omega_\beta}2\sum_{\{j,j'\}\subseteq G}R_{jj'}\pi_{\beta,j}\pi_{\beta,j'},
\end{equation}
with $\pi_\beta=\pi^{\rm in}_\beta=\pi^{\rm out}_\beta$ and $R_{jj'}$ the effective resistance for
the conductances $l_{\beta,j}q_{jj'}r_{\beta,j'}$. To see this, recall from the proof of
Proposition~\ref{prop:layer-constants}(iv) that $Q_G(z)$ is reversible with respect to $\nu_G$
for real $z\ge0$, so that its vectors $r_\beta(z)=r_\beta$ and $l_\beta(z)=l_\beta$ do not move
with $z$. By Proposition~\ref{prop:constant}(b) and the formula for $C_G(\beta,z)$ in
Section~\ref{sec:layer-setting}, for real $z\ge0$
\[
C_G(\beta,z)=\frac{(\mu^{(z)}\cdot r_\beta)(l_\beta\cdot e^{(z)})}{(4\pi)^{(k-2)/2}}\,
\frac{\sqrt{\tree(c(z))}}{\prod_{j\in G}\beta_j},\qquad
c_{jj'}(z)=l_{\beta,j}A(z)_{jj'}r_{\beta,j'} .
\]
Since $l_{\beta,j}c^-_j=a_\beta\pi_{\beta,j}$ and $c^+_{j'}r_{\beta,j'}=b_\beta\pi_{\beta,j'}$, we
have $c_{jj'}(z)=c_{jj'}(0)+z\omega_\beta\pi_{\beta,j}\pi_{\beta,j'}$. The derivatives at $z=0$ of
the logarithms of $\mu^{(z)}\cdot r_\beta$ and $l_\beta\cdot e^{(z)}$ are the first two terms
of~\eqref{eq:layer-c1}. The third term is the derivative of $\frac12\log\tree(c(z))$, by the rule
of Kirchhoff that $\partial\log\tree(c)/\partial c_e=R_e$, which says that $c_eR_e$ is the
weighted fraction of the spanning trees that contain the edge $e$. The function $C_G(\beta,z)$
of Theorem~\ref{thm:L-facet} below is holomorphic and agrees with this one for real $z\ge0$, so
its derivative at $0$ is also given by~\eqref{eq:layer-c1}.

\subsection{The facet theorem}\label{sec:layer-theorem}

The layer needs the local law of the weighted chain $Q_G(z)$, uniformly for complex $z$ near
$0$.

\begin{theorem}[the weighted facet chain]\label{thm:L-facet}
Let $G$ satisfy \textup{(H)} with $|G|=k-1\ge1$, let $i\notin G$ with $F=G\cup\{i\}$
irreducible, and let $\mathcal K\subset\Delta_G^\circ$ be compact. There are a number $z_0>0$,
not given explicitly, and a constant $C$, both depending only on $Q$, $\mu$, $i$ and
$\mathcal K$, and functions $J_G(\beta,z)$ and
$C_G(\beta,z)$ on $\mathcal K\times\{|z|<2z_0\}$, holomorphic in $z$, jointly continuous, real
for real $z$, equal to the rate and the local-law constant of Section~\ref{sec:layer-setting}
for real $0\le z<2z_0$, and with $C_G(\beta,0)=C_G(\beta)>0$, such that uniformly for
$\beta\in\mathcal K$, $S\ge1$ and complex $|z|\le z_0$
\begin{equation}\label{eq:L-facet}
f^{(z)}_G(S\beta)=S^{-(k-2)/2}\,C_G(\beta,z)\,e^{-S\,J_G(\beta,z)}\,(1+\varepsilon),\qquad|\varepsilon|\le C/S .
\end{equation}
For $|G|=1$ this holds with $\varepsilon=0$.
\end{theorem}

At $z=0$ the display is the local law~\eqref{eq:layer-local-G}, that is
Theorem~\ref{thm:local} for the face $G$. The proof inherits every constant from $z=0$ by
continuity, which is why $z_0$ is not explicit, and for larger $|z|$ the complex critical point
that it uses can cease to exist. For one or two states the word sum is a closed expression.

\begin{lemma}[one and two states]\label{lem:L-two}
Theorem~\ref{thm:L-facet} holds when $|G|\le2$, with the functions $J_G$ and $C_G$ given in the
proof.
\end{lemma}

\begin{proof}
For $G=\{j\}$ the proof of Proposition~\ref{prop:pair-face} gives
$f^{(z)}_G(s)=(\mu_j+z\mu_iq_{ij})(1+zq_{ji})e^{-(q_j-z\omega)s}$, which is~\eqref{eq:L-facet}
with $\varepsilon=0$, $J_G(\beta,z)=q_j-z\omega$ and
$C_G(\beta,z)=(\mu_j+z\mu_iq_{ij})(1+zq_{ji})$. Now let $G=\{1,2\}$, and write
$a_{12}(z)=q_{12}+zc^-_1c^+_2$ and $a_{21}(z)=q_{21}+zc^-_2c^+_1$ for
the entries of $A(z)$, $\varpi(z)=a_{12}(z)a_{21}(z)$, $p=\mu^{(z)}$ and $e=e^{(z)}$. A word in $G$
with no equal neighbours alternates. The words that start with $1$ are $1(21)^n$, with weight
$p_1\varpi^ne_1$ and $(n+1,n)$ letters, and $(12)^n$, with weight $p_1\varpi^{n-1}a_{12}e_2$ and
$(n,n)$ letters, for $n\ge1$. With $\xi=2\sqrt{\varpi s_1s_2}$ the power series of $I_1$ and $I_0$
give
\[
\sum_{n\ge1}\frac{\varpi^ns_1^ns_2^{n-1}}{n!\,(n-1)!}=\sqrt{\frac{\varpi s_1}{s_2}}\,I_1(\xi),\qquad
\sum_{n\ge1}\frac{(\varpi s_1s_2)^{n-1}}{(n-1)!^2}=I_0(\xi).
\]
We do the same for the words that start with $2$ and obtain
\[
f^{(z)}_G(s)=e^{-\langle d(z),s\rangle}\Bigl[\Bigl(p_1e_1\sqrt{\frac{\varpi s_1}{s_2}}+p_2e_2\sqrt{\frac{\varpi s_2}{s_1}}\Bigr)I_1(\xi)
+\bigl(p_1a_{12}e_2+p_2a_{21}e_1\bigr)I_0(\xi)\Bigr].
\]
Both sides are entire in $z$, so the branch of the square root does not matter. Since $G$ is
irreducible, $\varpi(0)=q_{12}q_{21}>0$. We choose $z_0$ so that $|\varpi(z)-\varpi(0)|\le\varpi(0)/2$
for $|z|\le2z_0$ and take the principal branch of $\sqrt\varpi$. Then
$\xi=2S\sqrt{\beta_1\beta_2\varpi}$ at $s=S\beta$ satisfies $|\arg\xi|\le\pi/12$ and $|\xi|\ge cS$,
uniformly on $\mathcal K$. The expansion $I_\nu(\xi)=e^\xi(2\pi\xi)^{-1/2}(1+O(1/|\xi|))$ holds
for $\nu=0,1$, uniformly in the sector $|\arg\xi|\le\pi/4$ and $|\xi|\ge1$~\cite[Eq.~10.40.1]{dlmf}.
Since $(2\pi\xi)^{-1/2}=S^{-1/2}(4\pi)^{-1/2}(\varpi\beta_1\beta_2)^{-1/4}$, this
gives~\eqref{eq:L-facet} with
\[
J_G(\beta,z)=\langle d(z),\beta\rangle-2\sqrt{\beta_1\beta_2\varpi},\qquad
C_G(\beta,z)=\frac{U+V}{\sqrt{4\pi}\,(\varpi\beta_1\beta_2)^{1/4}},
\]
where $U=p_1e_1\sqrt{\varpi\beta_1/\beta_2}+p_2e_2\sqrt{\varpi\beta_2/\beta_1}$ and
$V=p_1a_{12}e_2+p_2a_{21}e_1$, and with $\varepsilon=(U\varepsilon_1+V\varepsilon_0)/(U+V)$, where
$|\varepsilon_\nu|\le C/S$ are the errors of the two expansions. Both functions are holomorphic
in $z$, jointly continuous and real for real $z$. At $z=0$ we have $p=\mu_G$ and $e=\ones$, so
$U+V>0$ because $\mu(G)>0$. After shrinking $z_0$, $|U+V|$ is bounded below and $|U|+|V|$ above
on $\mathcal K\times\{|z|\le2z_0\}$, which gives $|\varepsilon|\le C/S$ once $cS\ge1$. For
$1\le S\le1/c$ the error $\varepsilon$ is continuous on a compact set, so it is bounded, and we
enlarge $C$. For real $z\ge0$ the
rate of the two-state matrix $Q_G(z)$ is
$\sup_u[\beta_1(d_1-a_{12}u_2/u_1)+\beta_2(d_2-a_{21}u_1/u_2)]=\langle d,\beta\rangle-2\sqrt{\beta_1\beta_2\varpi}$,
by the inequality of the arithmetic and geometric means. The constant agrees with that of
Section~\ref{sec:layer-setting}, because for real $z\ge0$ both are the limit of
$S^{1/2}e^{SJ_G(\beta,z)}f^{(z)}_G(S\beta)$. For the first this is the paragraph above. For the
second it is Lemma~\ref{lem:local-general} on $G$ for the weights $A(z)\ge0$, $\mu^{(z)}$ and
$e^{(z)}$, since $f^{(z)}_G(s)=e^{-\langle d(z),s\rangle}H_{A(z),\mu^{(z)},e^{(z)}}(s)$ in the
notation of Appendix~\ref{app:torus}.
\end{proof}

\emph{Facets with two or more states.} Let $|G|=k-1\ge2$. We number the states of $G$ as
$1,\dots,k-1$ and put $m=k-2$. For $z\in\mathbb C$ let $A'(z)$ be $A(z)$ without the row and
column of state $1$, and let $X_z$, $g_z$ and $b_z$ be the functions of
Lemma~\ref{lem:torus-complex} on $G$ for the weights $A(z)$, $\mu^{(z)}$ and $e^{(z)}$. Since
$c^-$, $c^+$ and $\mu_i$ are nonnegative, the weights $A(|z|)$, $\mu^{(|z|)}$ and $e^{(|z|)}$
dominate these in the sense of that lemma. For fixed $\beta\in\Delta_G^\circ$ and
$z'\in\mathbb C^m$ put
\[
\Phi_z(z')=\beta_1g_z(z')+\sum_{j\ge2}\beta_jz_j .
\]
Let $\theta(\beta)$ be the tilt of the face $G$ at $\beta$, let $\theta(\beta,z)$ be that of
$Q_G(z)$ for real $z\ge0$, and let $\mathcal H(\beta)$ be the Hessian $\mathcal H$ of the proof of
Theorem~\ref{thm:local} for the face $G$ at $\alpha=\beta$. On $\mathbb C^m$ we use the norm
$\lVert w\rVert=\max_j|w_j|$ with closed balls $\bar B(w,r)$, and $we^{\mathrm i\varphi}$ is the
vector with entries $w_je^{\mathrm i\varphi_j}$. The proof of Theorem~\ref{thm:local} for $G$
gives, for every $\beta\in\Delta_G^\circ$,
\begin{enumerate}[label=(P\arabic*)]
\item $\varrho(A'(0)-\operatorname{diag}\theta'(\beta))<0$ and $g_0(\theta'(\beta))=\theta_1(\beta)$,
by~\eqref{eq:A-radii},
\item $\nabla\Phi_0(\theta'(\beta))=0$, and $\nabla^2\Phi_0(\theta'(\beta))=\mathcal H(\beta)$ is
real, symmetric and positive definite,
\item $\operatorname{Re}\Phi_0(\theta'(\beta)e^{\mathrm i\varphi})\le\Phi_0(\theta'(\beta))
-\sum_{j\ge2}\beta_j\theta_j(\beta)(1-\cos\varphi_j)$ for $\varphi\in\mathbb T^m$,
\item $b_0(\theta'(\beta))>0$.
\end{enumerate}
Also $\theta$ and $\mathcal H$ are continuous, by Proposition~\ref{prop:constant}(a).

\begin{lemma}[the complex critical point]\label{lem:critical}
Let $|G|\ge2$ and let $\mathcal K\subset\Delta_G^\circ$ be compact. There are $z_1>0$,
$r_0\ge r_1>0$, $\sigma_0>0$ and $B_0$, which depend only on $Q$, $\mu$, $i$ and $\mathcal K$, such
that the following hold for every $\beta\in\mathcal K$.
\begin{enumerate}[(a)]
\item $r_0<\frac14\min_j\theta_j(\beta)$, and $\varrho(A'(|z|)-\operatorname{diag}s')<0$ for
$|z|\le z_1$ and real $s'$ with $\lVert s'-\theta'(\beta)\rVert\le3r_0$.
\item The functions $g_z(z')$ and $b_z(z')$ are holomorphic in $(z,z')$ and bounded in modulus by
$B_0$ on $\{|z|<z_1\}\times U_\beta$, where $U_\beta$ is the set of $z'$ with
$|z_j|>\theta_j(\beta)-3r_0$ for all $j$.
\item For $|z|<z_1$ there is exactly one $\zeta'(\beta,z)\in\bar B(\theta'(\beta),r_1)$ with
$\nabla\Phi_z(\zeta'(\beta,z))=0$. It is holomorphic in $z$ and jointly continuous in
$(\beta,z)$, $\zeta'(\beta,0)=\theta'(\beta)$, $\lVert\zeta'(\beta,z)-\theta'(\beta)\rVert\le B_0|z|$,
and $\zeta'(\beta,\bar z)=\overline{\zeta'(\beta,z)}$.
\item The matrix $\mathcal J(\beta,z)=Z_c\nabla^2\Phi_z(\zeta'(\beta,z))Z_c$, where
$Z_c=\operatorname{diag}\zeta'(\beta,z)$, is complex symmetric and jointly continuous, and
$\mathcal J(\beta,0)=\Theta'\mathcal H(\beta)\Theta'\ge\sigma_0I$, where
$\Theta'=\operatorname{diag}\theta'(\beta)$.
\item For real $0\le z<z_1$ we have $\zeta'(\beta,z)=\theta'(\beta,z)$,
$\Phi_z(\zeta'(\beta,z))=\langle\beta,\theta(\beta,z)\rangle$,
$\beta_1^2\det\mathcal J(\beta,z)=\prod_{j\ge2}\theta_j(\beta,z)^2\det\Sigma_\beta(z)$ and
$\beta_1b_z(\zeta'(\beta,z))=(\mu^{(z)}\cdot r_\beta(z))(l_\beta(z)\cdot e^{(z)})$.
\end{enumerate}
\end{lemma}

\begin{proof}
(a) The Perron root of a Metzler matrix is a continuous function of its entries. By (P1) the
function $(\beta,y,s')\mapsto\varrho(A'(y)-\operatorname{diag}s')$ is negative on the compact set
of the points $(\beta,0,\theta'(\beta))$ with $\beta\in\mathcal K$. So it is negative on a
neighbourhood of this set, and the neighbourhood contains the points with $\beta\in\mathcal K$,
$0\le y\le z_1$ and $\lVert s'-\theta'(\beta)\rVert\le3r_0$ when $z_1$ and $r_0$ are small. We also
take $r_0<\frac14\min_{\mathcal K}\min_j\theta_j$, which is positive.

(b) Put $s'_\beta=\theta'(\beta)-3r_0\ones>0$. By (a) we may use Lemma~\ref{lem:torus-complex}
with the dominating weights at $|z|$ and the radii $s'_\beta$. So $\mathcal Z'-A'(z)$ is invertible
on $\{|z|\le z_1\}\times U_\beta$, and there $|g_z|\le\bar g(s'_\beta)$ and
$|b_z|\le\bar b(s'_\beta)$ for the functions $\bar g$ and $\bar b$ of the weights at $|z|$. The
terms of their series do not decrease with $|z|$, so we may take the weights at $z_1$. Then the
bounds are continuous in $\beta$, since $(\operatorname{diag}s'-A'(z_1))^{-1}$ is continuous where
$\varrho(A'(z_1)-\operatorname{diag}s')<0$. So they are bounded on $\mathcal K$. The weights are
affine in $z$, so the last claim of Lemma~\ref{lem:torus-complex} gives the holomorphy.

(c) For $|z|\le z_1/2$ and $z'\in\bar B(\theta'(\beta),r_0)$ the polydisc of radius
$\min(z_1/2,r_0)$ about $(z,z')$ lies in $\{|z|<z_1\}\times U_\beta$, since its points have
$|z_j|\ge\theta_j-2r_0$. So by (b) and Cauchy's estimates, the derivatives of $g_z(z')$ and
$b_z(z')$ in $(z,z')$ up to order $4$ are bounded on $\{|z|\le z_1/2\}\times\bar B(\theta'(\beta),r_0)$
by a constant $B_1$ that does not depend on $\beta$. Let
$F(z,z')=\nabla\Phi_z(z')=\beta_1\nabla g_z(z')+\beta'$. By (P2), $F(0,\theta')=0$ and
$D_{z'}F(0,\theta')=\mathcal H(\beta)$, whose inverse has norm at most $B_2$ on $\mathcal K$ by
continuity. Put $T_z(z')=z'-\mathcal H(\beta)^{-1}F(z,z')$. For $|z|\le z_1/2$, $r\le r_0$ and
$z'\in\bar B(\theta',r)$ we get
\[
\lVert D_{z'}T_z(z')\rVert=\lVert\mathcal H^{-1}(D_{z'}F(0,\theta')-D_{z'}F(z,z'))\rVert\le B_3(|z|+r),\qquad
\lVert T_z(\theta')-\theta'\rVert\le B_3|z|,
\]
with $B_3$ depending only on $m$, $B_1$ and $B_2$. Choose $r_1\le r_0$ and $z_2\le z_1/2$ with
$B_3(z_2+r_1)\le\frac12$ and $B_3z_2\le\frac12r_1$. For $|z|\le z_2$ the map $T_z$ sends
$\bar B(\theta',r_1)$ into itself and is a contraction with constant $\frac12$ there. So $F(z,\cdot)$
has exactly one zero $\zeta'(\beta,z)$ in this ball, the fixed point of $T_z$, and
$\lVert\zeta'-\theta'\rVert\le2\lVert T_z(\theta')-\theta'\rVert\le2B_3|z|$. The iterates
$T_z^n(\theta')$ are holomorphic in $z$, continuous in $(\beta,z)$ and converge uniformly, which
gives the regularity of $\zeta'$, and $\zeta'(\beta,0)=\theta'$. The data $A_G$, $c^\pm$, $\mu$
and $\beta$ are real, so $F(\bar z,\bar w)=\overline{F(z,w)}$. Hence
$\overline{\zeta'(\beta,z)}$ is a zero of $F(\bar z,\cdot)$ in the ball, and it equals
$\zeta'(\beta,\bar z)$. We rename $z_2$ as $z_1$ and enlarge $B_0$ to at least $2B_3$.

(d) The Hessian $\nabla^2\Phi_z$ is symmetric and $Z_c$ is diagonal, so $\mathcal J$ is complex
symmetric, and it is continuous by (b) and (c). At $z=0$, $Z_c=\Theta'$ and (P2) gives
$\mathcal J(\beta,0)=\Theta'\mathcal H(\beta)\Theta'$, which is positive definite and continuous in
$\beta$. So its least eigenvalue has a positive minimum $\sigma_0$ on $\mathcal K$.

(e) Let $0\le z<z_1$ be real and write $\zeta'=\zeta'(\beta,z)$. By (c), $\zeta'$ is real, and it is
positive since $\lVert\zeta'-\theta'\rVert\le r_0<\frac14\min_j\theta_j$. The matrix
$A(z)=A(|z|)$ is nonnegative, and its graph contains that of $A_G$, so it is strongly connected.
By (a) and the first fact, $X_z(\zeta')\ge0$, and as in the proof of Theorem~\ref{thm:local} the
vector $\tilde r=(1,X_z(\zeta')b_1)\ge0$ lies in the kernel of $A(z)-\operatorname{diag}\tilde\theta$
with $\tilde\theta=(g_z(\zeta'),\zeta')$. By the second fact, $\tilde r>0$ and $0$ is the Perron
root. Let $\tilde l>0$ be the left Perron vector with $\tilde l\cdot\tilde r=1$, and
$\pi=\tilde l\circ\tilde r$. With $\tilde R=\operatorname{diag}\tilde r$, the matrix
$\tilde R^{-1}(A(z)-\operatorname{diag}\tilde\theta)\tilde R$ is an irreducible generator with
stationary law $\pi$, so Lemma~\ref{lem:perron} gives $\nabla\lambda(\tilde\theta)=-\pi$ for
$\lambda(y)=\varrho(A(z)-\operatorname{diag}y)$. For real $z'$ near $\zeta'$ the same argument
gives $\lambda(g_z(z'),z')=0$, and differentiating at $\zeta'$ gives
$\pi_1\partial_jg_z(\zeta')+\pi_j=0$. The critical equation $\beta_1\partial_jg_z(\zeta')+\beta_j=0$
then gives $\pi_j/\pi_1=\beta_j/\beta_1$ for $j\ge2$, and since $\pi$ and $\beta$ both sum to $1$,
$\pi=\beta$. So $\tilde l=\beta/\tilde r$, and $\tilde l^\top(A(z)-\operatorname{diag}\tilde\theta)=0$
says that the derivative of $u\mapsto\sum_{j\ne j'}\beta_jA(z)_{jj'}u_{j'}/u_j$ vanishes at
$\tilde r$, as before~\eqref{eq:stationary}. This function is convex in $\log u$, so $\tilde r$ is
a minimizer, and Proposition~\ref{prop:rate}(b) for $Q_G(z)$ shows that $\tilde r$ is a multiple
of $r_\beta(z)$. Hence $\tilde\theta=A(z)\tilde r/\tilde r=\theta(\beta,z)$, and
$\Phi_z(\zeta')=\beta_1\theta_1+\sum_{j\ge2}\beta_j\theta_j=\langle\beta,\theta(\beta,z)\rangle$.
The last two claims are the identities of Lemma~\ref{lem:local-general} on $G$ for $A(z)$,
$\mu^{(z)}$ and $e^{(z)}$ at $\alpha=\beta$, since $\mathcal J(\beta,z)$ is the Hessian
$\mathcal H$ of that lemma multiplied on both sides by $\operatorname{diag}\theta'(\beta,z)$.
\end{proof}

\begin{proof}[Proof of Theorem~\ref{thm:L-facet}]
For $|G|=1$ this is Lemma~\ref{lem:L-two}. Let $|G|\ge2$, and let $z_1$, $r_0$, $r_1$, $\sigma_0$
and $B_0$ be as in Lemma~\ref{lem:critical}. For $\beta\in\mathcal K$ and $|z|<z_1$ write
$\zeta'=\zeta'(\beta,z)$ and $\rho_j=|\zeta_j|$, and put
\[
\Psi(\beta,z,\varphi)=\Phi_z(\zeta'e^{\mathrm i\varphi}),\qquad
\chi(\beta,z,\varphi)=b_z(\zeta'e^{\mathrm i\varphi})\prod_{j\ge2}\zeta_je^{\mathrm i\varphi_j}.
\]

\emph{Step 1: the constants.} Since $\theta_j$ is real, $|\rho_j-\theta_j|\le|\zeta_j-\theta_j|\le r_1$.
So the points $\zeta'e^{\mathrm i\varphi}$ lie in $U_\beta$, and by Lemma~\ref{lem:critical}(b),
(c) the $\varphi$-derivatives of $\Psi$ and $\chi$ are jointly continuous on the compact set
$P_1\times\mathbb T^m$, where $P_1=\mathcal K\times\{|z|\le z_1/2\}$. Let $B_\Psi$ and $B_a$ bound
those of $\Psi$ up to order $4$ and those of $\chi$ up to order $2$. Since
$\partial/\partial\varphi_j=\mathrm iw_j\,\partial/\partial w_j$ at $w=\zeta'e^{\mathrm i\varphi}$
and $\nabla\Phi_z(\zeta')=0$, we have $\nabla_\varphi\Psi(\beta,z,0)=0$ and
$-\nabla^2_\varphi\Psi(\beta,z,0)=\mathcal J(\beta,z)$, as in the proof of Theorem~\ref{thm:local}.
Fix $\varepsilon\in(0,1)$ with $m^3B_\Psi\varepsilon\le\frac34\sigma_0$, and put
$2\eta=\beta_{\min}\theta_{\min}(1-\cos\varepsilon)$, where $\beta_{\min}$ and $\theta_{\min}$ are
the least entries of $\beta$ and $\theta(\beta)$ over $\mathcal K$. By (P3),
$\operatorname{Re}\Psi(\beta,0,\varphi)-\Psi(\beta,0,0)\le-2\eta$ for
$\lVert\varphi\rVert\ge\varepsilon$. By uniform continuity on $P_1\times\mathbb T^m$ we can choose
$z_0\in(0,z_1/4]$ such that for $\beta\in\mathcal K$ and $|z|\le2z_0$
\begin{enumerate}[(\roman*)]
\item $\operatorname{Re}\mathcal J(\beta,z)\ge\frac12\sigma_0I$, since
$\lVert\mathcal J(\beta,z)-\mathcal J(\beta,0)\rVert\le\frac12\sigma_0$ in the operator norm and
$\mathcal J(\beta,0)$ is real,
\item $|\chi(\beta,z,0)|\ge c_\chi$, where $2c_\chi>0$ is the minimum over $\mathcal K$ of
$\chi(\beta,0,0)=b_0(\theta')\prod_{j\ge2}\theta_j$, which is positive by (P4),
\item $\operatorname{Re}\Psi(\beta,z,\varphi)-\operatorname{Re}\Psi(\beta,z,0)\le-\eta$ for
$\lVert\varphi\rVert\ge\varepsilon$.
\end{enumerate}
The numbers $z_0$, $\varepsilon$, $\eta$ and $c_\chi$ depend only on $Q$, $\mu$, $i$ and
$\mathcal K$.

\emph{Step 2: the torus through the critical point.} Let $\beta\in\mathcal K$, $|z|\le z_0$ and
$S\ge1$. Since $\lVert\rho'-\theta'\rVert\le r_0$, Lemma~\ref{lem:critical}(a) shows that the
radii $\rho'$ are admissible in Lemma~\ref{lem:torus-complex} with the dominating weights at
$|z|$. With $t=S\beta$ it gives
\[
e^{S\langle d(z),\beta\rangle}f^{(z)}_G(S\beta)=\frac1{(2\pi)^m}\int_{\mathbb T^m}
\exp\bigl(S\Phi_z(w)\bigr)\,b_z(w)\prod_{j\ge2}w_j\,d\varphi,\qquad w_j=\rho_je^{\mathrm i\varphi_j}.
\]
The translation of $\varphi$ by the arguments of the $\zeta_j$ preserves the measure and turns $w$
into $\zeta'e^{\mathrm i\varphi}$. So the integral is $\int_{\mathbb T^m}e^{S\Psi}\chi\,d\varphi$.
No contour is moved. The torus representation holds at every admissible radius, and we use it at
the radii $\rho'$ of the complex critical point.

\emph{Step 3: the Laplace lemma.} By Step~1, Lemma~\ref{lem:laplace-complex} applies on
$\mathcal P=\mathcal K\times\{|z|\le z_0\}$ with $\sigma_{\min}=\frac12\sigma_0$, $\varepsilon$
and $\eta$. Since $\Psi(\beta,z,0)=\Phi_z(\zeta')$, it gives~\eqref{eq:L-facet} with
\[
J_G(\beta,z)=\langle d(z),\beta\rangle-\Phi_z(\zeta'),\qquad
C_G(\beta,z)=(2\pi)^{-m/2}\det\mathcal J(\beta,z)^{-1/2}\,b_z(\zeta')\prod_{j\ge2}\zeta_j,
\]
and with $|\varepsilon|\le C/(c_\chi S)$, where $C$ is the constant of
Lemma~\ref{lem:laplace-complex}.

\emph{Step 4: the functions.} We define $J_G$ and $C_G$ on $\mathcal K\times\{|z|<2z_0\}$ by these
formulas. Since $\zeta'$ is holomorphic in $z$ and $g_z(z')$ and $b_z(z')$ are holomorphic in
$(z,z')$, the numbers $\Phi_z(\zeta')$, $b_z(\zeta')$, $\prod_j\zeta_j$ and the entries of
$\mathcal J(\beta,z)$ are holomorphic in $z$ and jointly continuous. By (i) the matrix
$\mathcal J(\beta,z)$ lies in the set $\mathcal S_m$ of Appendix~\ref{app:laplace}, so
$\det\mathcal J(\beta,z)^{-1/2}$ is continuous, and as a continuous square root of a holomorphic
function without zeros it is holomorphic in $z$. For real $z$, $\zeta'$ is real by
Lemma~\ref{lem:critical}(c), so $J_G$ and $C_G$ are real, and $\mathcal J(\beta,z)$ is real and
positive definite, so the branch is the positive root. For real $0\le z<2z_0$,
Lemma~\ref{lem:critical}(e) gives $J_G(\beta,z)=\langle d(z)-\theta(\beta,z),\beta\rangle$, which is
the rate of $Q_G(z)$ by the identity $J=\langle q-\theta,\alpha\rangle$ after
Definition~\ref{def:tilt}, since the diagonal of $Q_G(z)$ is $-d(z)$. It also gives
$C_G(\beta,z)=(\mu^{(z)}\cdot r_\beta(z))(l_\beta(z)\cdot e^{(z)})/((2\pi)^{m/2}\sqrt{\det\Sigma_\beta(z)})$,
the constant of Section~\ref{sec:layer-setting}. At $z=0$ it is $C_G(\beta)>0$, by (H).
\end{proof}

The proof does not use the irreducibility of $F$, which only makes $c^-$ and $c^+$ nonzero.

We take the functions of Theorem~\ref{thm:L-facet} and define for $\beta\in\mathcal K$ and
$x,S>0$ the generalized kernel
\[
\mathbf K_\beta(x,S)=e^{-q_ix}\,\frac1{2\pi\mathrm i}\oint_{|z|=z_*}e^{x/z}\,\frac{C_G(\beta,z)}{C_G(\beta,0)}\,e^{S(J_G(\beta)-J_G(\beta,z))}\,\frac{dz}{z^2},
\]
for $z_*<2z_0$. The integral does not depend on the radius below $2z_0$.

\begin{theorem}[the facet layer]\label{thm:layer}
Let $F$ be irreducible with $k=|F|\ge2$, let $i\in F$, let $G=F\setminus\{i\}$ satisfy
\textup{(H)}, and let $\mathcal K\subset\Delta_G^\circ$ be compact. Then there
are $\delta>0$, $T_0$ and $C$ such that the following hold uniformly for $\beta\in\mathcal K$, $T\ge T_0$ and
$0<x\le\delta T$, with $S=T-x$ and $s=S\beta$.
\begin{enumerate}[(a)]
\item $f_F(s,x)=f_G(s)\,\mathbf K_\beta(x,S)\,(1+E)$ with $|E|\le C/T$.
\item $f_F(s,x)=f_G(s)\,\mathbf B_\beta(x,S)\,e^{E}$ with $|E|\le C\bigl(1/T+\sqrt{x/T}\,(1+x)\bigr)$.
\item If $Q$ is reversible on $F$, then $\gamma_\beta=q_i$, (b) holds with
$|E|\le C(1/T+\sqrt{x/T})$, and
\[
f_F(s,x)=f_G(s)\,\mathbf B_\beta(x,S)\Bigl(1+c_1(\beta)\,z_*\frac{I_0(\zeta)}{I_1(\zeta)}+E'\Bigr),\qquad|E'|\le C\,\frac{1+x}T .
\]
\end{enumerate}
For $k=2$, (a) holds with $E=0$, and Proposition~\ref{prop:pair-face} gives (b) for all
$x>0$ with an explicit bound.
\end{theorem}

When $Q$ is not reversible our bound gives the relative form $f_F=f_G\mathbf B_\beta(1+O(1/T+\sqrt{x/T}(1+x)))$ only for
$x\le CT^{1/3}$, and beyond that only the logarithmic form (b) and the kernel of (a). The
cause is the cubic term $S\,\partial_z^3J_G(\beta,0)z^3/6$ at $z\approx z_*$, of size
$x^{3/2}T^{-1/2}$, which the generalized kernel keeps and the Bessel kernel does not. For the
non-reversible cycle $q_{12}=q_{23}=q_{31}=2$, $q_{21}=q_{32}=q_{13}=1$ with $G=\{1,2\}$,
$\beta=(\frac12,\frac12)$ and $x=T/20$, the ratio $f_F/(f_G\mathbf B_\beta)$ is $19.6$, $265$ and
$48268$ at $T=800$, $1600$ and $3200$, so the relative form does fail there.

\begin{proof}[Proof of \textup{(a)}]
Theorem~\ref{thm:L-facet} gives~\eqref{eq:L-facet}, with $\varepsilon=0$ for $k=2$. Take
$\delta\le\frac12$, so $S\ge T/2$. Write $\Lambda(z)=S^{-(k-2)/2}C_G(\beta,z)e^{-SJ_G(\beta,z)}$.
By~\eqref{eq:L-facet}, $f_G(s)=\Lambda(0)(1+O(1/S))$, and by~\eqref{eq:layer-contour}, for every
$\rho\le z_0$,
\[
e^{q_ix}f_F(s,x)=\mathcal A+\mathcal D,\qquad \mathcal A=\frac1{2\pi\mathrm i}\oint_{|z|=\rho}e^{x/z}\Lambda(z)\frac{dz}{z^2},\qquad
|\mathcal D|\le\frac CS\,\mathcal A^{\rm abs},
\]
where $\mathcal A^{\rm abs}$ is the integral of $|e^{x/z}\Lambda(z)||z|^{-2}$ against $|dz|/(2\pi)$. By
definition $\mathbf K_\beta(x,S)=e^{-q_ix}\mathcal A/\Lambda(0)$. So it suffices to show that
$\mathcal A>0$ and $\mathcal A^{\rm abs}\le C'\mathcal A$. Put
\[
\Xi(z)=\frac{C_G(\beta,z)}{C_G(\beta,0)}\,e^{S(J_G(\beta)-J_G(\beta,z)-\omega_\beta z)},
\]
so that $\Lambda(z)/\Lambda(0)=e^{S\omega_\beta z}\Xi(z)$. By
Proposition~\ref{prop:layer-constants}(ii), the holomorphy of $C_G$, and Cauchy estimates on the
compact set $\mathcal K\times\{|z|\le\frac32z_0\}$, the function $R=\log\Xi$ is
holomorphic on $|z|\le z_0$ after shrinking $z_0$, and
\[
R(z)=c_1z-S\omega_\beta(\gamma_\beta-q_i)z^2+O(|z|^2+S|z|^3),\qquad |R''(z)|\le C(1+S).
\]
We take $\delta$ so small that $z_*\le\sqrt{2\delta/\omega_{\min}}\le z_0/2$, where
$\omega_{\min}=\min_{\mathcal K}\omega_\beta>0$.

\emph{Few visits, $\zeta\le Z_0$} for a fixed large $Z_0$. Take $\rho=\max(z_*,1/S)$, which is
at most $C/S$. On the circle $|x/z|\le Z_0/2$, $|S\omega_\beta z|\le C$ and
$|\Xi(z)-1|\le C/S$. Let $\mathcal A_0$ be $\mathcal A$ with $\Xi=1$. By
\eqref{eq:layer-bessel-moments} it equals $\Lambda(0)\sqrt{\omega_\beta S/x}\,I_1(\zeta)$, which
is at least $\Lambda(0)\omega_\beta S$, because
$\sqrt{\omega S/x}\,I_1(\zeta)=\omega S\sum_{n\ge0}(\zeta/2)^{2n}/(n!(n+1)!)$. Also
$|\mathcal A-\mathcal A_0|\le C\Lambda(0)\rho^{-1}/S\le C\Lambda(0)$ and
$\mathcal A^{\rm abs}\le C\Lambda(0)\rho^{-1}\le C\Lambda(0)S$. So
$\mathcal A\ge\Lambda(0)(\omega_\beta S-C)>0$ and $\mathcal A^{\rm abs}\le C'\mathcal A$ for $S$ large.

\emph{Many visits, $\zeta\ge Z_0$.} Let $\lambda(z)=x/z+S\omega_\beta z+R(z)$, so that
$e^{x/z}\Lambda(z)=\Lambda(0)e^{\lambda(z)}$. It is real on the real axis, and
$\lambda'(z)=-x/z^2+S\omega_\beta+O(1+S|z|)$. By the implicit function theorem it has a real
critical point $\rho_s=z_*(1+O(\sqrt{x/S}+1/\zeta))$, and we take $\rho=\rho_s$. For
$z=\rho e^{\mathrm i\phi}$,
\[
\operatorname{Re}\lambda(z)-\lambda(\rho)=(x/\rho+S\omega_\beta\rho)(\cos\phi-1)+\operatorname{Re}[R(z)-R(\rho)] .
\]
Since $R$ is real on the real axis, $\operatorname{Re}[R'(\rho)(z-\rho)]=R'(\rho)\rho(\cos\phi-1)$,
and $|R''|\le C(1+S)$ bounds the rest by $2C(1+S)\rho^2(1-\cos\phi)$. Also
$|R'(\rho)|\rho\le C(\rho+S\rho^2)$. Now $S\rho^2\le Cx$, $x/\rho+S\omega_\beta\rho\ge\zeta$,
and $x/\zeta\le C\sqrt\delta$. So for $\delta$ small
$\operatorname{Re}\lambda(z)-\lambda(\rho)\le-\frac\zeta2(1-\cos\phi)$, and
\[
\mathcal A^{\rm abs}\le\Lambda(0)\rho^{-1}e^{\lambda(\rho)}\frac1{2\pi}\int_{-\pi}^\pi e^{-\frac\zeta2(1-\cos\phi)}d\phi
\le C\Lambda(0)\rho^{-1}e^{\lambda(\rho)}\zeta^{-1/2}.
\]
For the lower bound we write
$\mathcal A=\Lambda(0)\rho^{-1}\frac1{2\pi}\int e^{\zeta\ell(\phi)}e^{-\mathrm i\phi}d\phi$ with
$\ell(\phi)=\lambda(\rho e^{\mathrm i\phi})/\zeta$. Here $\ell(0)$ is real, $\ell'(0)=0$,
$-\ell''(0)=\rho^2\lambda''(\rho)/\zeta$ is real and lies in $[\frac12,2]$, the derivatives of
$\ell$ up to order four are bounded uniformly in all parameters, and the gap above holds. The
parameters $(\beta,x,S)$ do not range over a compact set, so we use the uniform form of
Lemma~\ref{lem:laplace} in Remark~\ref{rem:laplace-uniform}, with $\sigma_{\min}=\frac12$ and the
gap $\frac12(1-\cos\varepsilon)$. With $m=1$, $\Psi=\ell$, $a=e^{-\mathrm i\phi}$ and $\zeta$ in place
of $T$ it gives
$\mathcal A=\Lambda(0)\rho^{-1}e^{\lambda(\rho)}(2\pi\zeta(-\ell''(0)))^{-1/2}(1+O(1/\zeta))$,
which is positive and at least $c\,\mathcal A^{\rm abs}$.
\end{proof}

\begin{proof}[Proof of \textup{(b)} and \textup{(c)}]
We use two facts on Bessel functions besides~\eqref{eq:layer-bessel-moments}. For integers
$n\ge1$, $I_n(\zeta)=\frac1\pi\int_0^\pi e^{\zeta\cos t}\cos(nt)\,dt$~\cite[Eq.~10.32.3]{dlmf}.
So $I_n\le I_1$ for all $\zeta>0$, and $|1-\cos nt-(1-\cos t)|\le n^2t^2/2$ and a Laplace
estimate give $|I_1(\zeta)-I_n(\zeta)|\le Cn^2I_1(\zeta)/\zeta$ for $\zeta\ge1$. Also
$I_0/I_1\le1+2/\zeta$, as in the proof of Proposition~\ref{prop:pair-face}.

\emph{The reversible case.} By Proposition~\ref{prop:layer-constants}(iv) and holomorphy,
$J_G(\beta,z)=J_G(\beta)-z\omega_\beta$ on the disc, so
$\Xi(z)=C_G(\beta,z)/C_G(\beta,0)=1+c_1z+\sum_{n\ge2}c_nz^n$ with $|c_n|\le Cz_0^{-n}$. By
\eqref{eq:layer-bessel-moments} the term $c_nz^n$ contributes $c_nz_*^nI_{n-1}(\zeta)/I_1(\zeta)$
relative to the main term $z_*^{-1}I_1(\zeta)$. For $n\ge2$ this is at most $|c_n|z_*^n$, and the
sum over $n\ge2$ is $O(z_*^2)=O(x/T)$, since $z_*\le z_0/2$. The term $n=1$ is
$c_1z_*I_0(\zeta)/I_1(\zeta)=O(z_*+1/(\omega_\beta S))=O(\sqrt{x/T}+1/T)$. With (a) this proves
both claims of (c), for all $0<x\le\delta T$.

\emph{The general case, $x\le1$.} Write $\Xi(z)=e^{-S\omega_\beta(\gamma_\beta-q_i)z^2}(1+c_1z+R_2(z))$
with $|R_2(z)|\le C(|z|^2+S|z|^3)$. On $|z|=z_*$ the Gaussian factor has modulus at most
$e^{(\gamma_\beta-q_i)x}\le C$. We expand it in powers of $z^2$ and use
\eqref{eq:layer-bessel-moments}. With $y=(\gamma_\beta-q_i)x$ the term $1$ gives the relative
factor $\sum_{m\ge0}(-y)^mI_{2m-1}(\zeta)/(m!\,I_1(\zeta))$. For $\zeta\ge1$ it differs from
$e^{-y}$ by at most $C\sum_{m\ge1}y^mm^2/(m!\,\zeta)\le Cx/\zeta=O(\sqrt{x/S})$. For $\zeta<1$
we have $x<1/(4\omega_\beta S)$, and $|I_{2m-1}/I_1-1|\le1$ bounds the difference by
$e^y-1=O(1/S)$. In the same way the term $c_1z$ gives $O(z_*(1+1/\zeta))=O(\sqrt{x/S}+1/S)$.
The term $R_2$ is bounded through the modulus of the integrand, and
$\frac1{2\pi}\int e^{\zeta\cos\phi}d\phi=I_0(\zeta)$ gives
$O\bigl((z_*^2+Sz_*^3)I_0(\zeta)/I_1(\zeta)\bigr)=O(1/S+\sqrt{x/S})$. So
$\mathbf K_\beta=\mathbf B_\beta(1+O(1/T+\sqrt{x/T}))$ for $x\le1$.

\emph{The general case, $1\le x\le\delta T$.} Here $\zeta\ge2\sqrt{\omega_{\min}S}\ge Z_0$. We
use the evaluation of $\mathcal A$ in the proof of (a), and the same evaluation with
$\Xi=1$, whose integral is exactly $\Lambda(0)z_*^{-1}I_1(\zeta)$. For $\Xi=1$ the
critical point is $z_*$, the critical value is $\zeta$, and $-\ell''(0)=1$. The ratio of the two
evaluations is $e^{\lambda(\rho_s)-\zeta}(1+O(\sqrt{x/S}+1/\zeta))$, since
$\rho_s/z_*-1$ and the relative change of $-\ell''(0)$ are $O(\sqrt{x/S}+1/\zeta)$. Now
$x/z+S\omega_\beta z$ takes the value $\zeta$ at its critical point $z_*$, and the shift from
$z_*$ to $\rho_s$ changes $\lambda$ only by
$O(\lambda''(\rho_s)(\rho_s-z_*)^2)=O(\zeta x/S+1/\zeta)$. So
\[
\lambda(\rho_s)-\zeta=R(z_*)+O(\zeta x/S+1/\zeta)=-(\gamma_\beta-q_i)x+O\bigl(\sqrt{x/S}\,(1+x)\bigr),
\]
because $c_1z_*=O(\sqrt{x/S})$, $S\omega_\beta z_*^2=x$, $Sz_*^3=O(x^{3/2}S^{-1/2})$,
$\zeta x/S=O(x^{3/2}S^{-1/2})$ and $1/\zeta\le C\sqrt{x/S}$ for $x\ge1$. Hence
$\mathbf K_\beta=\mathbf B_\beta\exp(O(\sqrt{x/T}(1+x)))$, and with (a) this proves (b).
\end{proof}

At $x=\alpha_iT$ with $\alpha_i$ small but fixed, the saddle point
of $\mathbf K_\beta$ reproduces the exponent $-TJ_F(\alpha)$ exactly, by the Legendre identity of
Proposition~\ref{prop:layer-constants}(i). So part (a) and Theorem~\ref{thm:local} overlap, and
together they cover the face up to its facet.

\subsection{The lattice and the near-vertex case}\label{sec:layer-vertex}

Take $Q=J-dI$, the uniform start and $\ell=1$ in Section~\ref{sec:vertex}, with $F=\{1,2\}$,
$G=\{1\}$, $i=2$ and $n=(b,a)$. Then $q_{12}=q_{21}=1$, $p=1-hq_1$ and $u=h/p$, and
Lemma~\ref{lem:layer-exact-discrete} gives exactly
\[
P_N(b,a)=\sum_{K\ge1}\binom{a-1}{K-1}(1-hq_2)^{a-K}\,[z^K]\,(\mu_1+zh\mu_2)(1+zh)(1-hq_1+zh^2)^{b-1}.
\]
Its main part is $\mu_1p^{N-1}\sum_K\binom{a-1}{K-1}\binom{b-1}Ku^{2K}$, which becomes
$B_a(bu^2)$ of~\eqref{eq:refresh-polynomial} when $\binom{b-1}K$ is replaced by $b^K/K!$. So for
$\ell=1$ the near-vertex law is the discrete pair-face case, with $\omega=1$, $S=hb$, $x=ha$ and
$\gamma=q_2$. The parameter of Theorem~\ref{thm:vertex-uniform-root} is $w=bu^2=\omega Sh/p^2$,
which is $\omega Sh$ up to the factor $p^{-2}=1+O(h)$. There $w$ is of order one, Lemma~\ref{lem:uniform-laguerre-bessel} pins $B_a(w)/(w\mathcal I(aw))$ between
$1-w/2$ and $1$, and the continuum layer does not apply. When $T^2/N\to0$ the parameter $w$
tends to $0$, and the same lemma shows that the lattice adds no separate regime.

\emph{An open problem.} A discrete form of Theorem~\ref{thm:layer}, with the Laguerre
polynomial $B_{n_i}$ in place of the Bessel kernel, would need a discrete form of
Theorem~\ref{thm:L-facet} and the saddle bounds of the proof above on the lattice, and we do not
state it.

The scripts of Section~\ref{sec:verify} check the statements of this section on three chains
with three states, including the cycle above, for which $\omega_\beta=5.1820$ and
$\gamma_\beta=3.4179$ against $q_3=3$. The generalized kernel matches $f_F/f_G$ to $O(1/T)$ for
$x$ from $0.01/T$ to $T/4$ on all three. They also check Theorem~\ref{thm:L-facet} on two facets
with four states, a facet of a non-reversible face with five states and a path whose jump chain
is periodic. There the exact word sum and the prediction differ by a relative error $\varepsilon$
with $S|\varepsilon|$ settling at a constant, for instance $0.438$ and $0.545$ at $z=0$, with the
ratio of the values at $S=32$ and $S=16$ equal to $1.002$ and $1.007$. No statement is deduced
from these numbers.
